% GENERATED FILE. Edit canonical lessons, then run npm run generate:books.
% canonical-source-hash: fnv1a64:cad6aaed9687d7e7
% canonical-lessons: TA-C75-polar, TA-C75-answering, TA-W30-read-sariyaa, TA-C75-permission

\chapter{Yes or No}
\label{ch:ta-yes-or-no}
\hlchaptermodality{75}

\begin{chapteropening}
\textbf{By the end of this chapter:} \emph{I can turn any sentence I know into a yes-or-no question, answer one, and ask permission.}

Complete TA-C75-permission and ask to be let in, then grant the same request.
\end{chapteropening}

\section[-ā]{-\ta{ஆ} — the question you can build yourself}

\label{lesson:TA-C75-polar}



\begin{quote}
Read \textbf{\ta{பாலும்}}. Then say: neither milk nor tea.
\end{quote}

\begin{grammarlens}[title={Grammar lens: one letter on the end}]
Every question you have been given so far came ready-made: \ta{எப்படி}, \ta{என்ன}, \ta{யார்}, \ta{எங்கே}, \ta{எவ்வளவு}. Each one asks about a thing you name in advance.

\textbf{-\ta{ஆ}} (\emph{-ā}) asks about the whole sentence, and it is yours to attach to anything.

Take a statement. Put \textbf{-\ta{ஆ}} on its last word. It is now a question that wants \emph{yes} or \emph{no}.

\noindent
\begin{tabularx}{\linewidth}{@{}>{\raggedright\arraybackslash}X>{\raggedright\arraybackslash}X@{}}
\toprule
\textbf{statement} & \textbf{question} \\
\midrule
\textbf{\ta{நீங்கள்} \ta{நலம்}} — you are well & \textbf{\ta{நீங்கள்} \ta{நலமா}?} — are you well? \\
\textbf{\ta{இது} \ta{பால்}} — this is milk & \textbf{\ta{இது} \ta{பாலா}?} — is this milk? \\
\textbf{\ta{சரி}} — okay & \textbf{\ta{சரியா}?} — okay? \\
\bottomrule
\end{tabularx}

Word order does not move. Nothing is added at the front. The ending goes on the \textbf{last} word, and the last word is the one being asked about.
\end{grammarlens}

\begin{grammarlens}[title={What you've built}]
Every sentence you already own can now be a question.
\end{grammarlens}

\subsection*{Your turn}
\begin{itemize}
\raggedright
  \item \emph{Say it:} \emph{sariyā?}
  \item \emph{Say it:} \emph{nīṅgaḷ nalamā?}
  \item \emph{Turn:} say \emph{idu pāl}, then ask \emph{idu pālā?}
  \item \emph{Recall:} read \textbf{\ta{பாலும்}}, then say \emph{pālum tēnīrum illai}, then ask \emph{sariyā?}
  \item \emph{Return to:} say \emph{āṇḍu}, say \emph{tākam} and read \textbf{\ta{ஓ}} — three distances back — then ask about one of them with -\ta{ஆ}
\end{itemize}

\begin{tcolorbox}[breakable,colback=teal!4,colframe=teal!35!black,title={Before you move on}]
Which word in the sentence carries \textbf{-\ta{ஆ}}? (\textbf{The last one — and that is the word the question is about.})
\end{tcolorbox}
\section[ām / illai]{\ta{ஆம்}, \ta{இல்லை} — and the answer Tamil likes better}

\label{lesson:TA-C75-answering}



\begin{quote}
Ask, in Tamil: is this milk? Then say: neither milk nor tea.
\end{quote}

\begin{grammarlens}[title={Grammar lens: two answers, and a third}]
You have had \textbf{\ta{ஆம்}} and \textbf{\ta{இல்லை}} since chapter one. Now there is something to answer.

\begin{quote}
\textbf{\ta{இது} \ta{பாலா}?} — \emph{idu pālā?}
\end{quote}

\begin{quote}
\textbf{\ta{ஆம்}.} / \textbf{\ta{இல்லை}.}
\end{quote}

Both work, and both are correct. But listen to how Tamil speakers actually reply, because they often reach for neither: they \textbf{repeat the word that was asked about}.

\begin{quote}
\textbf{\ta{இது} \ta{பாலா}?} — \textbf{\ta{பால்}.}
\end{quote}

\begin{quote}
\textbf{\ta{சரியா}?} — \textbf{\ta{சரி}.}
\end{quote}

The echo is the everyday answer. \textbf{\ta{ஆம்}} is a shade more formal, and it is the one to use when the question was long and you want to be plain.

\textbf{\ta{இல்லை}} has no echo version. To deny, you say \textbf{\ta{இல்லை}}.
\end{grammarlens}

\begin{grammarlens}[title={What you've built}]
You can ask a yes-or-no question and answer one, in the two shapes a speaker will hear.
\end{grammarlens}

\subsection*{Your turn}
\begin{itemize}
\raggedright
  \item \emph{Answer:} \emph{idu pālā?} — with the echo
  \item \emph{Answer:} \emph{idu pālā?} — with \ta{இல்லை}
  \item \emph{Ask and answer:} \emph{sariyā?} — then \emph{sari}
  \item \emph{Recall:} say \emph{pālum tēnīrum illai}, then ask \emph{nīṅgaḷ nalamā?}, then answer \emph{nalam}
  \item \emph{Return to:} say \emph{-um … -um}, \emph{tūkkam} and \emph{ōḍu} — three distances back — then ask about one of them with -\ta{ஆ}
\end{itemize}

\begin{tcolorbox}[breakable,colback=teal!4,colframe=teal!35!black,title={Before you move on}]
Somebody asks \textbf{\ta{சரியா}?} and you agree. What is the everyday reply? (\textbf{\ta{சரி}} — the echo, not \ta{ஆம்}.)
\end{tcolorbox}
\section[sariyā]{\ta{சரியா} — a question you can see}

\label{lesson:TA-W30-read-sariyaa}



\begin{quote}
Before the page: ask \emph{sariyā?}, then give the echo answer.
\end{quote}

\subsection*{Script you'll notice: \ta{சரியா} — the ending is a sign, not a letter}
\textbf{\ta{சரி}} was the third word in this book, and you wrote it in chapter nineteen. The question ending adds one shape to it.

\noindent
\begin{tabularx}{\linewidth}{@{}>{\raggedright\arraybackslash}X>{\raggedright\arraybackslash}X>{\raggedright\arraybackslash}X@{}}
\toprule
\textbf{piece} & \textbf{} & \textbf{} \\
\midrule
\textbf{\ta{ச}} & \emph{ca} & the letter for a whole family of sounds \\
\textbf{\ta{ரி}} & \emph{ri} & \ta{ர} with the \textbf{\ta{ி}} sign \\
\textbf{\ta{யா}} & \emph{yā} & \ta{ய} with the long \textbf{\ta{ா}} sign \\
\bottomrule
\end{tabularx}

\begin{quote}
\textbf{\ta{ச} · \ta{ரி} · \ta{யா} $\to$ \ta{சரியா}}
\end{quote}

A \textbf{\ta{ய}} appeared that was not in \textbf{\ta{சரி}}. That is the same slide you saw in \ta{மேசையும்}: \ta{சரி} ends in a vowel, so a light \textbf{\ta{ய்}} carries the ending on.

And the ending itself is the sign \textbf{\ta{ா}}, sitting after \ta{ய} the way it sits after \ta{ப} in \textbf{\ta{பா}}. A question mark closes the line on the page, but the question is already in that sign.

\subsection*{Your turn}
\begin{itemize}
\raggedright
  \item \emph{Read it:} \textbf{\ta{ச}} — then \textbf{\ta{ரி}} — then \textbf{\ta{யா}}
  \item \emph{Read it:} \textbf{\ta{சரி}}, then \textbf{\ta{சரியா}} — a statement, then a question
  \item \emph{Point to:} the shape that does the asking
  \item \emph{Recall:} ask \emph{sariyā?}, then answer \emph{sari}, then read \textbf{\ta{சரியா}}
  \item \emph{Return to:} say \emph{-um}, \emph{kāyccal} and \emph{kaḻuvu} — three distances back — then ask about one of them with -\ta{ஆ}
\end{itemize}

\begin{tcolorbox}[breakable,colback=teal!4,colframe=teal!35!black,title={Before you move on}]
Read \textbf{\ta{சரியா}}. Which single shape turned it into a question? (\textbf{The long \ta{ா} on the end.})
\end{tcolorbox}
\section[nāṉ uḷḷē varalāmā?]{\ta{நான்} \ta{உள்ளே} \ta{வரலாமா}? — may I come in?}

\label{lesson:TA-C75-permission}



\begin{quote}
Read \textbf{\ta{சரியா}}, then answer it.
\end{quote}

\subsection*{You'll want to know: \ta{வரலாமா}?}
\textbf{\ta{நான்} \ta{உள்ளே} \ta{வரலாமா}?} (\emph{nāṉ uḷḷē varalāmā?}) — \textquotedblleft{}may I come in?\textquotedblright{}

Three pieces you have, and one you half-have. \ta{நான்} is \textquotedblleft{}I\textquotedblright{}. \ta{உள்ளே} is \textquotedblleft{}inside\textquotedblright{}, from the six directions. \ta{வா} is \textquotedblleft{}come\textquotedblright{}.

The half-known piece is \textbf{-\ta{லாம்}}. It has been in your goodbye since chapter four: \textbf{\ta{நாளை} \ta{பார்க்கலாம்}} — \textquotedblleft{}we may see each other tomorrow\textquotedblright{}. On a verb, \textbf{-\ta{லாம்}} means \emph{it is allowed}.

\ta{வா} becomes \textbf{\ta{வரலாம்}} — \textquotedblleft{}may come\textquotedblright{}. Then chapter seventy-five's ending arrives:

\begin{quote}
\textbf{\ta{வரலாம்}} $\to$ \textbf{\ta{வரலாமா}?} — \textquotedblleft{}may I come?\textquotedblright{}
\end{quote}

Permission is a yes-or-no question, so it takes the yes-or-no ending. Nothing new to learn about the asking.

To grant it, you have both answers already:

\begin{quote}
\textbf{\ta{ஆம்}, \ta{வரலாம்}.} or, at a doorway, plainly: \textbf{\ta{சரி}.}
\end{quote}

\begin{grammarlens}[title={What you've built}]
A doorway you can stand in and be let through.
\end{grammarlens}

\subsection*{Your turn}
\begin{itemize}
\raggedright
  \item \emph{Say it:} \emph{nāṉ uḷḷē varalāmā?}
  \item \emph{Grant:} \emph{ām, varalām}
  \item \emph{Grant again:} \emph{sari}
  \item \emph{Recall:} answer \emph{sariyā?}, then read \textbf{\ta{சரியா}}, then ask \emph{varalāmā?}
  \item \emph{Return to:} read \textbf{\ta{பாலும்}}, say \emph{vali} and say \emph{viḷaiyāḍu} — three distances back — then ask about one of them with -\ta{ஆ}
\end{itemize}

\begin{tcolorbox}[breakable,colback=teal!4,colframe=teal!35!black,title={Before you move on}]
You are at a door and want to be let in. Say it. (\textbf{\ta{நான்} \ta{உள்ளே} \ta{வரலாமா}?})
\end{tcolorbox}
